\documentclass[11pt,a4paper]{article}
\usepackage[T1]{fontenc}\usepackage{lmodern,amsmath,amssymb,amsthm,mathtools,booktabs,microtype}
\usepackage[margin=26mm]{geometry}\usepackage[colorlinks=true,linkcolor=blue,urlcolor=blue,citecolor=blue]{hyperref}
\newtheorem{theorem}{Theorem}[section]\newtheorem{lemma}[theorem]{Lemma}\newtheorem{proposition}[theorem]{Proposition}\newtheorem{corollary}[theorem]{Corollary}
\newcommand{\Q}{\mathbb Q}\newcommand{\Z}{\mathbb Z}\newcommand{\F}{\mathbb F}\DeclareMathOperator{\Tr}{Tr}
\setlength{\emergencystretch}{2em}
\title{An inert-prime obstruction to reflection matching of pole inputs}
\author{}\date{15 November 2026}
\begin{document}\maketitle
\begin{abstract}
Reflection makes the two eight-contact detector inputs equal and preserves
the elementary height comparison. Nevertheless it is not a complete
restriction of the current conditional quartic for any sufficiently large
odd coefficient shift. A calculation at the inert prime 11 proves that,
for every such shift, some integer between 1 and 124 has no witness.
All four possible torsion translates of the recovered elliptic point are
excluded. A separate calculation at 19 gives infinitely many shifts at which
the fixed input 2 fails. These results concern this explicit restriction,
not all possible constructions with matching pole square classes.
\end{abstract}

\section{The source and the proposed restriction}
Retain the complete source and arithmetic premises of the 79-variable
construction \cite{Q79}. Put
\[
 K=\Q(\rho),\quad \rho^2=2,\quad d=3-\rho,\quad D=d^2,
 \quad E:y^2=x^3-Dx,
\]
\[
 P=\left((57-28\rho)/8,(-660+417\rho)/32\right),\qquad
 2E(K)=\Z P.
\]
The source uses, for a fixed odd coefficient $B\ge3$,
\[
 \beta=X^2+B,\quad T=\Tr x(Q),\quad a=\beta^4+T^4,
 \quad t_0=\lambda,\quad t_1=\lambda+a.
\]
Its recovered finite point $Q$ has index $\nu(Q)\in\tfrac12\Z$, defined
by $2Q=[2\nu(Q)]P$. The retained source soundness theorem says that,
for every sufficiently large fixed odd $B$, every source solution satisfies
$\nu(Q)=\beta$. We use precisely that implication under its stated premises.

The contact function and the proposed restriction are
\begin{equation}\label{eq:reflection}
 F(t)=\prod_{r\in\{1,2,3,6\}}((t-1)^2-r^2),\qquad
 \lambda=1-a/2,\quad t_1=2-\lambda.
\end{equation}
Reflection makes $F(t_0)=F(t_1)$ and $(t_0-1)^2=(t_1-1)^2$.
Since the two detectors have identical field tests, their nonzero physical
pole inputs agree, not just their square classes. For the rational logarithmic
height $h(u/v)=\log\max(|u|,|v|)$ in lowest terms, the affine relation gives
\[
 h(a)\le h(\lambda)+\log4,\qquad
 h(\lambda)\le h(a)+\log4.
\]
The failure below is therefore not a failure of square-class matching or
of this elementary height comparison.

\begin{theorem}\label{thm:uniform}
Under the source soundness premises, for every sufficiently large odd
$B$ there is an integer $X$ with $1\le X\le124$ for which the source
restricted by~\eqref{eq:reflection} has no rational witness.
In particular no sufficiently large member of this restricted family
defines all of $\Z$.
\end{theorem}

\section{A local obstruction at a small inert prime}
We will use $p=11$ and $p=19$. Both are inert in $K$, both give good
reduction of $E$, and the eight contacts are distinct modulo either prime.
Write $v$ for the unique valuation, normalized by $v(p)=1$.
For a point of $E(K)$ every finite nonzero reciprocal abscissa $f=1/x$
has even valuation. If $v(x)<0$, the cubic term in the curve equation
proves that $v(x)$ is even; if $v(x)>0$, its linear term does the same,
because $D$ is a unit. Extend $f(O)=0$ with infinite valuation.

\begin{lemma}\label{lem:inert}
No rational source solution has a positive odd contact valuation at
$p=11$ or $p=19$.
\end{lemma}
\begin{proof}
The source's exceptional-case proof recovers $H,H-P$ before any division
by contact weights. Suppose a contact has positive odd valuation $\eta$.
Then $t-1$ is a unit and $\eta=v(F(t))$. Write
\[
 w=F(t)/(t-1)^6,\qquad q=w^{-1},\qquad v(q)=-\eta.
\]
The literal source quadrics give the two physical factors
\[
 W_1=1+qC^2+q^2 f(H)^2+q^4 f(U)^4,
\]
\[
 W_2=1+qI^2+q^2\sigma f(H-P)^2+q^4\sigma f(V)^4.
\]
Nonzero nonconstant summands have valuations, respectively, odd,
congruent to 2 modulo 4, and congruent to 4 modulo 8. The three classes
are disjoint and none contains zero. Thus each factor is nonzero, has
valuation at most zero, and its valuation is the unique minimum of its
nonzero summands and zero. A vanishing test is simply omitted.

For $W=W_1W_2$ the source trace rows give the nonzero pole input
\[
 b=\frac{F(t)}{(t-1)^{10}\Tr(W^{-1})^2}.
\]
For clarity, if the stored product is $\Omega$, the row
$\Omega_0^2-2\Omega_0\eta_0-2\Omega_1^2=0$ excludes $\Omega_0=0$
when $\Omega\ne0$, and then
$\eta_0=\operatorname{N}(\Omega)/\Tr(\Omega)\ne0$.
Hence the displayed trace really is nonzero. Since $p$ is inert,
$v(\Tr(W^{-1}))\ge-v(W)$, giving
\[
                       v(b)\le\eta+2v(W).
\]
This valuation is odd. The congruent-number pole equations exclude negative
odd valuations at every odd prime \cite{Pole}, so $v(b)>0$ and
$v(W)>-\eta/2$. Each factor valuation is at least $v(W)$.
The two terms with reciprocal reference abscissas consequently imply
\[
 v(f(H))>3\eta/4>0,\qquad v(f(H-P))>3\eta/4>0.
\]
Both reference points reduce to the identity at the same place, so their
difference $P$ would do so. The finite-field computations below show that
it does not. No formal-logarithm truncation or large-prime height radical
is used in this argument. It therefore legitimately applies at 11 and 19,
even though the earlier height estimates omitted these primes.
\end{proof}

\section{The complete obstruction at 11}
In $\F_{121}=\F_{11}[\rho]/(\rho^2-2)$ direct arithmetic gives
\begin{equation}\label{eq:p11}
 P=(3+2\rho,\rho),\quad
 2P=(8+9\rho,8),\quad 3P=(8+9\rho,3)=-2P.
\end{equation}
Thus $P$ has order five. The endomorphism $(x,y)\mapsto(-x,iy)$ with
$i=4\rho$, $i^2=-1$, sends $P$ to $2P$. It permutes $E[2]$.
The traces of the four torsion translates of $P$, in the order
$O,(0,0),(d,0),(-d,0)$, are
\[
                             6,7,9,9\pmod {11}.
\]
For example, their exact rational traces before reduction are
\[
 \frac{57}{4},\quad-\frac{4656}{1681},\quad
 \frac{5670}{289},\quad-\frac{338}{147}.
\]
Every nonzero multiple of the reduced $P$ is $\pm P$ or $\pm iP$.
Since $x(-Q)=x(Q)$ and $x(iQ)=-x(Q)$, it follows that, whenever
$n\not\equiv0\pmod5$, all four translates of $[n]P$ are affine and
\begin{equation}\label{eq:fourthmultiset}
 \bigl\{(\Tr x([n]P+\tau))^4:\tau\in E[2]\bigr\}
                         =\{9,3,5,5\}\quad\hbox{in }\F_{11}.
\end{equation}
Here equality includes multiplicities.

Let $\mathcal C=\{\pm2,\pm4,\pm6,\pm12\}$, the possible signed doubled
contact distances. Its eight residues modulo 11 are distinct. If
$n\equiv\pm1$ or $\pm3\pmod {11}$, then $n^4$ is 1 or 4, and
\eqref{eq:fourthmultiset} shows that
\begin{equation}\label{eq:allcontacts}
 n^4+(\Tr x([n]P+\tau))^4\in\mathcal C\pmod {11}
                 \quad\hbox{for all four }\tau.
\end{equation}

\begin{lemma}\label{lem:smallinput}
For every integer $B$, there exists $x_0\in\{1,2,3,9,10,14\}$ such that
\[
 11\nmid x_0,\qquad B+x_0^2\not\equiv0\pmod5,
 \qquad B+x_0^2\equiv\pm1\hbox{ or }\pm3\pmod {11}.
\]
\end{lemma}
\begin{proof}
For $B\bmod11=0,1,\ldots,10$, initially choose respectively
\[
                       x_0=1,3,1,3,2,3,2,1,2,1,2.
\]
The corresponding residues $B+x_0^2$ are
$1,10,3,1,8,3,10,8,1,10,3$, which have the required form.
If $x_0=1$ and $B\equiv4\pmod5$, replace $x_0$ by 10.
If $x_0=2$ and $B\equiv1\pmod5$, replace it by 9.
If $x_0=3$ and $B\equiv1\pmod5$, replace it by 14.
These replacements preserve its square modulo 11 and remove precisely
the possible zero residue modulo five. All choices are nonzero modulo 11.
\end{proof}

\begin{proof}[Proof of Theorem~\ref{thm:uniform}]
Fix $B$ beyond the source soundness threshold and choose $x_0$ by
Lemma~\ref{lem:smallinput}. For $j=-2,-1,0,1,2$ put
\[
 X_j=x_0+55j,\qquad n_j=B+X_j^2.
\]
The congruences in that lemma remain true for $n_j$.
Since $[5]P$ reduces to the identity and multiplication by 11 kills the
first formal-group quotient, $[55]P$ is the identity modulo $11^2$.
But $n_j-n_0=110x_0j+3025j^2$ is divisible by 55. Thus, for each fixed
$\tau\in E[2]$, the affine trace
$T_\tau(j)=\Tr x([n_j]P+\tau)$ is constant modulo 121.
All these points stay affine since $n_j\not\equiv0\pmod5$.

Put $a_\tau(j)=n_j^4+T_\tau(j)^4$. For each $\tau$ there is a unique
$c_\tau\in\mathcal C$ congruent to $a_\tau(j)$ modulo 11, by
\eqref{eq:allcontacts}. Moreover
\[
 a_\tau(j)-a_\tau(0)\equiv440n_0^3x_0j\pmod {121}.
\]
After division by 11 the coefficient is
$7n_0^3x_0\not\equiv0\pmod {11}$. Each translate therefore forbids
at most one residue class of $j$ modulo 11 for which
$a_\tau(j)-c_\tau$ is divisible by 121. Four translates forbid at most
four of the five trial values. For a surviving $j$ one has
\[
                    v_{11}(a_\tau(j)-c_\tau)=1
                         \qquad\hbox{for every }\tau\in E[2].
\]
Take $X=|X_j|$, so $1\le X\le124$. A putative rational witness must,
by source soundness, have $\nu(Q)=\beta=n_j$. Thus
$Q=[n_j]P+\tau$ for some two-torsion $\tau$.
The reflection restriction gives $(t_0-1)^2=a_\tau(j)^2/4$.
Exactly one contact has valuation one at 11: the other factors are units,
as the contacts are distinct and $c_\tau$ is a unit. This contradicts
Lemma~\ref{lem:inert}, for every possible torsion translate.
\end{proof}

The argument is uniform in $B$, but uses the existing ineffective soundness
threshold rather than computing it. The finite checker additionally
reconstructs all 44 nonidentity residue classes of $nP$ modulo 121,
verifies them by independent additions, and checks a witness input for
every $B\bmod605$. The modulus is $\operatorname{lcm}(55,121)$.
Those finite checks confirm the arithmetic just used, not the external
soundness theorem.

\section{A fixed-input obstruction at 19}
A complementary calculation gives a failure at the same input for an
infinite set of shifts. In $\F_{19^2}$ one has
\[
 P=(6\rho,15-3\rho),\quad 2P=(8+\rho,8-\rho),\quad
 4P=(11-\rho,13+5\rho),\quad 5P=(16+\rho,0)=(-d,0).
\]
Thus $P$ has order ten. The four translated trace residues are $0,2,2,3$.
Set $n_t=121+190t$. Their point traces are constant modulo $19^2$,
because $[190]P$ is the identity modulo $19^2$. For the four translates,
$n_t^4+T_\tau^4$ has residue $7,4,4,12$ modulo 19, respectively.
Choose signed distances $c_\tau=-12,4,4,12$.
The first-order variation is
\[
 (n_t^4+T_\tau^4)-(n_0^4+T_\tau^4)\equiv38t\pmod {361}.
\]
Each translate excludes at most one residue $t$ modulo 19 from valuation
one. Hence at least fifteen of the nineteen progressions
\[
 B\equiv117+190t\pmod {3610},\qquad 0\le t<19,
\]
have no witness for $X=2$ once $B$ exceeds the retained soundness threshold.
At least one of $t=0,1,2,3,4$ works. This gives a positive relative density
at least $15/1805=3/361$ among odd shifts for failure at this single input.
The proof is the same four-translate argument as above.

\section{Consequences and limitations}
Theorem~\ref{thm:uniform} rules out reflection as a complete restriction
for every sufficiently large fixed odd shift, even though it makes the
pole inputs identical and retains the height comparison. The obstruction
is stronger than the finiteness of a fixed-scale square-class curve.
It identifies an actual local incompatibility and a uniformly bounded
set of integer inputs detecting it.

This is not an impossibility theorem for matching first square classes by
other means. A new complete construction must control their local support,
not only their equality. Neither this paper nor its finite certificates
establishes a smaller Diophantine definition. The separate shifted-Lucas
quartic improvement retains independent poles.

\begin{thebibliography}{9}
\bibitem{Q79} Project manuscript,
\emph{A 79-variable quartic under rational $abc$},
\nolinkurl{papers/quartic_existence/quartic_abc79_20261114_v27.tex}.
\bibitem{Pole} Project manuscript,
\emph{Pole equations from simultaneous congruent-number twists},
\nolinkurl{papers/pole_parameters/congruent_twist_poles_20261114_v27.tex}.
\bibitem{Chart} Project manuscript,
\nolinkurl{papers/elliptic_charts/compressed_reference_20261114_v27.tex}.
\end{thebibliography}
\end{document}
